% Methodology
\section{The SPOQ Methodology}
\label{sec:methodology}

SPOQ orchestrates multi-agent software development through a four-stage pipeline:
Epic Planning, Epic Validation, Agent Execution, and Agent Validation. We present
the formal framework, key algorithms, and design principles underlying each stage.

\subsection{Design Principles}
\label{subsec:principles}

SPOQ is built on five design principles that distinguish it from prior multi-agent approaches:

\begin{principle}[Atomic Task Boundaries]
\label{prin:atomic}
Each task constitutes 1--4 hours of focused work with one clear deliverable.
Tasks are self-contained: a worker agent can complete the task without
coordinating with other agents during execution.
\end{principle}

\begin{principle}[Explicit Dependencies]
\label{prin:dependencies}
Task dependencies form a directed acyclic graph (DAG), a diagram where arrows
indicate prerequisite relationships, and no task can indirectly depend on itself.
All dependencies are declared upfront, enabling static analysis of parallelism
and identification of the \textit{critical path} (the longest chain of dependent
tasks that determines minimum project duration).
\end{principle}

\begin{principle}[Quality Gates Before and After]
\label{prin:gates}
Validation occurs both before execution (planning quality) and after execution
(code quality). Early validation catches expensive mistakes when they are cheapest
to fix.
\end{principle}

\begin{principle}[Human-AI Collaboration]
\label{prin:human}
A human specialist participates in task decomposition, validates plans, and can
be consulted by agents. The human is a high-value agent, not an external observer.
\end{principle}

\begin{principle}[Cost-Optimized Agent Selection]
\label{prin:cost}
Different roles require different capability-cost tradeoffs. High-capability
agents (Opus) handle complex tasks; balanced agents (Sonnet) handle quality
review; fast-cheap agents (Haiku) handle triage.
\end{principle}

\subsection{Four-Stage Pipeline}
\label{subsec:pipeline}

Figure~\ref{fig:pipeline} illustrates the SPOQ pipeline.

\begin{figure}[htbp]
\centering
\begin{tikzpicture}[
    node distance=0.8cm,
    stage/.style={rectangle, draw=spoqnavy, fill=spoqnavy!10, thick,
                  minimum width=3cm, minimum height=1cm, align=center},
    gate/.style={diamond, draw=spoqemerald, fill=spoqemerald!10, thick,
                 aspect=2, inner sep=2pt, align=center},
    arrow/.style={-Stealth, thick},
    label/.style={font=\small\itshape, text=spoqgray}
]

% Stages
\node[stage] (plan) {Epic Planning};
\node[gate, right=of plan] (eval1) {Epic\\Valid?};
\node[stage, right=of eval1] (exec) {Agent\\Execution};
\node[gate, right=of exec] (eval2) {Code\\Valid?};
\node[stage, right=of eval2] (done) {Complete};

% Arrows
\draw[arrow] (plan) -- (eval1);
\draw[arrow] (eval1) -- node[above, label] {$\geq$95} (exec);
\draw[arrow] (exec) -- (eval2);
\draw[arrow] (eval2) -- node[above, label] {$\geq$95} (done);

% Feedback loops
\draw[arrow, spoqburgundy] (eval1.south) -- ++(0,-0.5) -| node[below, label, pos=0.25] {revise} (plan.south);
\draw[arrow, spoqburgundy] (eval2.south) -- ++(0,-0.8) -| node[below, label, pos=0.25] {rework} (exec.south);

% Human integration
\node[above=0.3cm of plan, font=\small, text=spoqbrass] {Human + AI};
\node[above=0.3cm of eval1, font=\small, text=spoqbrass] {Human + AI};

\end{tikzpicture}
\caption{SPOQ four-stage pipeline with dual validation gates. Human specialist
participates in planning and plan validation. Failed validations trigger
revision/rework loops.}
\label{fig:pipeline}
\end{figure}

\begin{definition}[Epic]
\label{def:epic}
An \textit{epic} $E = (G, T, D, S)$ consists of:
\begin{itemize}
    \item $G$: Goal statement describing the desired outcome
    \item $T = \{t_1, t_2, \ldots, t_n\}$: Set of atomic tasks
    \item $D \subseteq T \times T$: Dependency relation where $(t_i, t_j) \in D$
          means $t_i$ must complete before $t_j$ can begin
    \item $S = \{s_1, s_2, \ldots, s_m\}$: Success criteria for the epic
\end{itemize}
\end{definition}

\begin{definition}[Task]
\label{def:task}
A \textit{task} $t = (id, desc, deps, files, criteria, est)$ consists of:
\begin{itemize}
    \item $id$: Unique identifier within the epic
    \item $desc$: Implementation description with steps
    \item $deps \subseteq T$: Set of prerequisite tasks (work that must finish first)
    \item $files$: List of files to be modified
    \item $criteria$: Acceptance criteria for task completion
    \item $est = (o, r, p)$: Three-point estimate: best-case ($o$), expected ($r$),
          and worst-case ($p$) durations, borrowed from PERT project management
\end{itemize}
\end{definition}

\subsection{Wave-Based Topological Dispatch}
\label{subsec:waves}

Given a task dependency graph, SPOQ computes \textit{waves}, groups of tasks
that can execute in parallel because they have no dependencies on each other.
The algorithm uses \textit{topological sorting}, a standard graph algorithm
that orders nodes (tasks) such that all dependencies appear before the nodes
that require them.

\begin{definition}[Wave Assignment]
\label{def:wave}
A \textit{wave assignment} $W: T \rightarrow \mathbb{N}$ maps each task to a
non-negative integer such that:
\begin{equation}
\forall (t_i, t_j) \in D: W(t_i) < W(t_j)
\end{equation}
Tasks in the same wave have no dependencies between them and can execute
concurrently.
\end{definition}

Algorithm~\ref{alg:waves} presents the wave computation procedure.

\begin{algorithm}[htbp]
\caption{Wave Computation via Topological Sort}
\label{alg:waves}
\begin{algorithmic}[1]
\Require Task set $T$, dependency relation $D$
\Ensure Wave assignment $W: T \rightarrow \mathbb{N}$

\State $\text{indegree}[t] \gets |\{t' : (t', t) \in D\}|$ for all $t \in T$
\State $W[t] \gets \bot$ for all $t \in T$
\State $w \gets 0$

\While{$\exists t \in T : W[t] = \bot$}
    \State $\text{ready} \gets \{t \in T : W[t] = \bot \land \text{indegree}[t] = 0\}$
    \For{$t \in \text{ready}$}
        \State $W[t] \gets w$
    \EndFor
    \For{$t \in \text{ready}$}
        \For{$t' : (t, t') \in D$}
            \State $\text{indegree}[t'] \gets \text{indegree}[t'] - 1$
        \EndFor
    \EndFor
    \State $w \gets w + 1$
\EndWhile

\State \Return $W$
\end{algorithmic}
\end{algorithm}

\begin{theorem}[Parallelism Bound]
\label{thm:parallelism}
Let $W^* = \max_{t \in T} W(t)$ be the number of waves. The wall-clock execution
time is bounded below by:
\begin{equation}
T_{\text{wall}} \geq \sum_{w=0}^{W^*} \max_{t: W(t)=w} \text{duration}(t)
\end{equation}
This bound is achieved when sufficient agents are available to execute all
tasks in each wave simultaneously.
\end{theorem}

\paragraph{Critical Path Analysis.}
The \textit{critical path} is the longest chain of dependent tasks through the
graph, weighted by durations. Even with unlimited parallel resources, the project
cannot complete faster than this path; it represents the irreducible sequential
bottleneck. Identifying critical path tasks reveals where delays have the greatest
impact. The theoretical minimum execution time is:

\begin{equation}
T_{\text{critical}} = \max_{\text{path } P} \sum_{t \in P} \text{duration}(t)
\end{equation}

SPOQ reports the \textit{speedup factor} $\sigma = T_{\text{sequential}} / T_{\text{critical}}$,
comparing sequential execution time (all tasks one-by-one) to the parallelized
critical path time. A speedup of 5.3x, for example, means the work completes
in roughly one-fifth the time it would take serially.

\subsection{Three-Tier Agent Hierarchy}
\label{subsec:agent-hierarchy}

SPOQ employs three agent tiers, each optimized for its role:

\begin{table}[H]
\centering
\caption{SPOQ Agent Hierarchy}
\label{tab:agents}
\begin{tabular}{llll}
\toprule
\textbf{Tier} & \textbf{Model} & \textbf{Role} & \textbf{Tradeoff} \\
\midrule
Worker & Opus & Task execution & High capability, high cost \\
Reviewer & Sonnet & Quality assurance & Balanced capability/cost \\
Investigator & Haiku & Build failure triage & Low cost, fast response \\
\bottomrule
\end{tabular}
\end{table}

\paragraph{Opus Workers.}
For each task $t$ in the current wave, SPOQ spawns an Opus agent with:
(1) the task specification, (2) relevant epic context, (3) completed dependency
outputs, and (4) prior QA feedback if this is a rework attempt.

\paragraph{Sonnet Reviewers.}
Each completed task undergoes quality review by a Sonnet agent, which scores
the work against 10 code quality metrics (Section~\ref{sec:validation}).
Tasks scoring below threshold are queued for rework with specific remediation
guidance.

\paragraph{Haiku Investigators.}
When the build fails after a wave, a Haiku agent analyzes the error output
to determine which tasks likely caused the failure. This triage is fast and
inexpensive, allowing rapid identification of problematic tasks without
engaging the full QA process.

\paragraph{Platform Independence.}
While our reference implementation uses Anthropic's Claude model family,
the SPOQ methodology is inherently platform-agnostic. The three-tier
hierarchy represents an abstract capability mapping: any sufficiently capable
model family can populate the Worker (high capability, high cost),
Investigator (low cost, fast response), and Reviewer (balanced capability/cost)
tiers. Table~\ref{tab:tier-mapping} illustrates potential mappings across providers.

\begin{table}[htbp]
\centering
\caption{Capability Tier Mapping Across LLM Providers}
\label{tab:tier-mapping}
\footnotesize
\begin{tabular}{lllll}
\toprule
\textbf{Tier} & \textbf{Role} & \textbf{Claude} & \textbf{OpenAI} & \textbf{Gemini} \\
\midrule
Worker & Task execution & Opus & GPT-4 & Ultra/Pro \\
Reviewer & Quality assurance & Sonnet & GPT-4-turbo & Pro \\
Investigator & Build triage & Haiku & GPT-3.5 & Flash \\
\bottomrule
\end{tabular}
\end{table}

The core methodological contributions, including wave-based topological dispatch,
dual validation gates, explicit dependency DAGs, and Human-as-an-Agent
integration, require no vendor-specific features and transfer directly
to alternative platforms. Organizations can implement SPOQ using their
preferred LLM provider by calibrating capability tiers to their available
models and adjusting cost thresholds accordingly.

\newpage
\subsection{Human-as-an-Agent (HaaA) Integration}
\label{subsec:haaa}

SPOQ treats the human specialist as a high-value agent integrated into the
orchestration loop, not an external supervisor. Rather than simply monitoring
AI agents from the sideline, the human actively participates, contributing
expertise where it matters most and receiving assistance in return.

\begin{definition}[Human-as-an-Agent]
\label{def:haaa}
The \textit{Human-as-an-Agent} (HaaA) paradigm defines bidirectional integration,
meaning communication flows both ways between human and system:
\begin{enumerate}
    \item \textbf{Human$\rightarrow$System:} The human participates in epic
    planning, validates epics before execution, and can intervene during execution.
    \item \textbf{System$\rightarrow$Human:} Agents can request human assistance
    when facing ambiguity, blocked progress, or decisions beyond their scope.
\end{enumerate}
\end{definition}

This bidirectional model enables \textit{quality amplification}: the human's
judgment improves task decomposition quality (reducing downstream rework),
while agents' execution scales the human's productivity.

\paragraph{Task Decomposition.}
The human specialist drafts epics using hierarchical task decomposition,
assisted by LLM suggestions. The human ensures:
\begin{itemize}
    \item Tasks are appropriately scoped (1--4 hours)
    \item Dependencies are correctly identified
    \item Acceptance criteria are verifiable
    \item Potential risks are mitigated
\end{itemize}

\paragraph{Validation Participation.}
Before execution, the human reviews the epic alongside the automated
validation. The human can approve, request revisions, or override
automated assessments with justification.

\paragraph{Consultation Requests.}
During execution, agents may encounter situations requiring human judgment:
ambiguous requirements, conflicting dependencies, or decisions with
significant implications. Agents can pause and request human input rather
than proceeding with potentially incorrect assumptions.

\begin{figure}[htbp]
\centering
\begin{tikzpicture}[
    node distance=1.8cm and 2.5cm,
    entity/.style={rectangle, draw=spoqnavy, fill=spoqnavy!10, thick,
                   minimum width=2.5cm, minimum height=0.8cm, align=center},
    human/.style={rectangle, draw=spoqbrass, fill=spoqbrass!10, thick,
                  minimum width=2.5cm, minimum height=0.8cm, align=center},
    arrow/.style={-Stealth, thick},
    bidir/.style={Stealth-Stealth, thick, spoqemerald}
]

% Entities
\node[human] (human) {Human Specialist};
\node[entity, below left=of human] (plan) {Epic Planning};
\node[entity, below=of human] (valid) {Validation};
\node[entity, below right=of human] (exec) {Execution};

% Bidirectional arrows with separated labels
\draw[bidir] (human) to[bend right=15] node[left, font=\small] {draft/refine} (plan);
\draw[bidir] (human) -- node[fill=white, font=\small, inner sep=2pt] {review/approve} (valid);
\draw[bidir] (human) to[bend left=15] node[right, font=\small] {consult/intervene} (exec);

% System flow
\draw[arrow, dashed] (plan) -- (valid);
\draw[arrow, dashed] (valid) -- (exec);

\end{tikzpicture}
\caption{Human-as-an-Agent integration. Bidirectional arrows indicate two-way
communication: human contributes to each stage and can be consulted by the system.}
\label{fig:haaa}
\end{figure}
